\chapter{Interview preparation}

\section{The 60-second pitch}
\sayit{I built an edge-first triage pipeline for gigapixel H\&E slides. A whole slide is a 100k by 200k pyramid, and
about 70\% of it is glass. On a thumbnail, a two-component Gaussian mixture on saturation and optical density finds the
tissue with Dice 0.84, against 0.53 for Otsu. On each 20$\times$ tile I compute 23 physically meaningful features -- focus,
stain, ink, folds, noise -- and score them with a GMM density learned only from clean tissue. That gives 0.94 AUROC for
artefacts, catches artefact types it never saw, with a 45 KB model that scores in microseconds and explains every
rejection. A pixel-level GMM in optical-density space also separates nuclei, stroma, fluid, blood and fat, and those
fractions alone give 0.93 AUROC for tumour vs normal. Only the good tiles go to the segmentation network, which in my
second project is a residual U-Net with a dilated bottleneck, patch stitching and test-time augmentation.}

\section{Likely questions with model answers}
\begin{description}
\item[Why GMM and not a CNN for triage?] Density estimation without labels; novel artefacts (64--99\% recall on unseen types
vs 7--64\% for supervised); 45\,KB; interpretable; calibrated threshold = chosen false-rejection rate. A CNN is better
when labelled examples of a failure mode exist (0.976) -- I would add one as a second head.
\item[Why GMM and not k-means?] k-means = GMM with identical spherical covariances and hard assignment; real clusters
in OD space are elongated and of very different size. Measured tissue Dice 0.84 vs 0.14; k-means never isolates fluid.
\item[How did you choose $K$?] Held-out-slide log-likelihood (BIC over-estimates because tiles of a slide are correlated).
K=1 is already strong -- with few slides the simplest density generalises best.
\item[How do you set the threshold?] 95th percentile of scores on clean tiles of \emph{slides not used for fitting}
(leave-one-slide-out). Same-slide calibration rejected 43\% of tumour tiles.
\item[How do you separate fluid?] Beer--Lambert $\to$ optical density $\to$ colour deconvolution into H and E; pixel GMM on
[H, E, OD, $a^*$, local std]; the pale, eosin-dominant, low-OD component is fluid (proteinaceous fluid/plasma/oedema);
nuclei = haematoxylin-dominant; blood = high $a^*$. Soft posteriors give a confidence map. Validation: 79\% agreement across
slides, k-means cannot separate it, fractions predict tumour at 0.93.
\item[What is the VQ-VAE for, then?] Compression/tokenisation of tiles and a tumour-novelty signal (0.75). Its MSE is a bad QC
score (0.38) because blurred and empty tiles are easy to reconstruct -- I verified that and changed the design.
\item[Correlation vs covariance?] Covariance has units and scale; correlation is covariance of standardised variables in
$[-1,1]$. Our features have $|\rho|>0.9$ pairs and a covariance condition number $10^{13}$ -- hence regularised covariances.
\item[KL vs JS?] KL asymmetric, unbounded, undefined with zero bins; JS symmetric, bounded $[0,1]$; we used JS to measure
which feature an artefact shifts and the stain shift between slides.
\item[Precision vs recall in QC?] Recall = fraction of bad tiles caught; precision = fraction of rejections that were bad;
specificity = good tissue kept. We fix specificity (95\% nominal) and maximise recall; and we check tumour is not rejected.
\item[IoU vs Dice?] Dice $=2IoU/(1+IoU)$, both overlap measures; Dice weights the intersection more; mIoU averages over
classes so small classes (cars, tumour) count equally.
\item[Gaussian vs motion blur? How detect?] Gaussian = isotropic low-pass (defocus); motion = directional box PSF.
Detect via Laplacian variance/Tenengrad/FFT high-frequency ratio; motion blur shows directional spectral zeros.
\item[Explain Canny.] Gaussian smoothing, Sobel gradients, non-maximum suppression, hysteresis with two thresholds.
\item[Data augmentation in pathology?] $D_4$ rotations/flips, HED stain jitter, brightness/contrast, blur/noise, elastic;
never augment away the label; TTA at inference.
\item[Supervised vs unsupervised?] Labels vs none; QC is novelty detection (unsupervised), diagnosis is supervised or weakly
supervised (MIL with slide labels).
\item[U-Net and ResNet?] Encoder--decoder with skip concatenations; ResNet adds identity shortcuts for gradient flow. Our
satellite ``residual'' is concatenation -- I fixed the terminology.
\item[VAE vs VQ-VAE?] Continuous Gaussian latent with KL to $\mathcal{N}(0,I)$ (ELBO, reparameterisation) vs discrete codebook
with commitment loss and straight-through gradients; VQ-VAE avoids posterior collapse and gives tokens.
\item[YOLOv8 vs YOLO11?] Both anchor-free one-stage; v8 C2f blocks + decoupled head; 11 adds C3k2 and C2PSA attention,
fewer parameters for higher mAP; -seg heads predict prototype masks. Use for countable objects (mitoses, nuclei, cars).
\item[Transformers in pathology?] ViT/Swin backbones, foundation models (UNI, Virchow, CONCH), MIL aggregators
(ABMIL, CLAM, TransMIL). Expensive; triage reduces the number of tiles they embed.
\item[Why ONNX Runtime?] Framework-independent graph, optimised CPU/GPU kernels, INT8; but measure: INT8 gave 3.7$\times$
smaller models and no speed-up on our CPU.
\item[How do you avoid leakage?] Split by slide/patient, never by tile; calibrate on unseen slides.
\item[Biggest limitation?] Synthetic artefacts (no public artefact labels); 7 slides; tissue GT itself from a CNN; fluid has
no pixel ground truth. Next step: a small annotated artefact set from the lab's own scanners.
\end{description}

\section{Do and don't}
\begin{itemize}
\item DO quote only numbers from this report and say how they were measured (slides, split, CPU).
\item DO mention the bugs you found and fixed -- it shows ownership.
\item DON'T claim SOTA. Classical Potsdam mIoU is 46--48 vs 87 for deep SOTA; that is the point: GMMs for colour/density
problems, CNNs for shape/semantics.
\item DON'T call concatenation a residual connection, or say scatter\_nd if the code loops.
\end{itemize}
